% ============================================================================
%  Reporte tecnico: simulacion del TFIM con qubits logicos vs fisicos
%  17 de octubre de 2026. Todo lo que no gasta creditos esta terminado;
%  falta solo correr en el emulador H2 (HQC), que es un cambio en config.py.
%
%  Compila con pdfLaTeX (motor por defecto). Autocontenido: solo paquetes estandar.
%  Figuras: las genera `scripts/plot_q3.py` en figures/. Si falta una, se imprime un
%  recuadro y el documento compila igual.
% ============================================================================
\documentclass[11pt,a4paper]{article}

\usepackage[utf8]{inputenc}
\usepackage[T1]{fontenc}
\usepackage{lmodern}
\usepackage[spanish,es-noindentfirst,es-noshorthands,es-tabla]{babel}
\usepackage[margin=2.5cm]{geometry}
\usepackage{amsmath}
\usepackage{graphicx}
\usepackage{booktabs}
\usepackage{enumitem}
\usepackage{hyperref}
\usepackage{xurl}
\hypersetup{colorlinks=true, linkcolor=black, urlcolor=blue, citecolor=black}
\setlength{\emergencystretch}{3em}

\newcommand{\code}[1]{\nolinkurl{#1}}
\newcommand{\pred}{\textsf{\small[predicci\'on cl\'asica, ruido asumido; no es H2]}}

\newcommand{\figura}[2]{%
  \IfFileExists{../../figures/#1}{\includegraphics[width=#2]{../../figures/#1}}{%
  \IfFileExists{../figures/#1}{\includegraphics[width=#2]{../figures/#1}}{%
  \IfFileExists{figures/#1}{\includegraphics[width=#2]{figures/#1}}{%
  \fbox{\parbox{0.8\linewidth}{\centering\small Falta la figura \code{#1}.\\
  Generarla con \code{python -I scripts/plot_q3.py}, o subirla a una carpeta \code{figures/}.}}}}}}

\title{\textbf{Simulaci\'on del TFIM con qubits l\'ogicos vs f\'isicos}\\[0.3em]
\large Reporte t\'ecnico para el equipo}
\author{}
\date{17 de octubre de 2026 \quad\textbullet\quad listo para el emulador del H2}

\begin{document}
\maketitle

\begin{abstract}
\noindent Simulaci\'on del TFIM que compara el c\'alculo de observables f\'isicos con qubits l\'ogicos
(codificados) frente a qubits f\'isicos, y que estudia la frustraci\'on, dirigida al hardware
Quantinuum H2. La \emph{ruta de env\'io a H2} est\'a construida y probada con las
capas gratuitas de Quantinuum; todo lo que no cuesta cr\'editos est\'a terminado. Lo \'unico que falta
es correr en el emulador H2 (gasta cr\'editos HQC), y eso es cambiar un par de l\'ineas de
\texttt{config.py} y dar una aprobaci\'on expl\'icita por corrida. Todas las curvas con ruido son
predicciones cl\'asicas bajo un modelo de ruido asumido, nunca resultados del H2.
\end{abstract}

\tableofcontents

% ============================================================================
\section{Introducci\'on}

La parte cl\'asica del proyecto est\'a terminada y el camino hacia el H2 est\'a construido y probado en las
capas gratuitas de Quantinuum (el verificador de sintaxis y el emulador sin ruido alojado en Nexus), sin
gastar un solo cr\'edito.

\textbf{C\'omo se hizo.} Partí de todo lo que habíamos decidido, de la discusión previa, la teoría que conocíamos y lo que le pusimos en la aplicación de Quantinuum, y usé eso para delimitarlo el trabajo usar Claude Code. Yo superviso las decisiones y trabajé con Claude Code. Al principio Kiro
(otro agente) ayud\'o a planificar, investigar y revisar; desde el 8 de octubre Kiro est\'a fuera, as\'i
que ahora es solo Claude Code planifica, construye y, para
revisar, lanza revisores independientes ``a ciegas'' (no ven el razonamiento interno, solo buscan errores y corroboran). Lo hice a consciencia, pero por favor rev\'isenlo.

Cada problema que encontré en el software de Quantinuum (Guppy/Selene, qnexus, pytket-quantinuum, PECOS, TKET) está en \code{docs/quantinuum_issues.md} con su reproducción~\cite{qilog}.

\subsection{Resumen del proyecto}
Pongo esto para estar en la misma página, o por si hay algún malentendido de mi parte.

El objetivo final es comparar, de frente, el c\'alculo de los mismos observables f\'isicos del TFIM
usando qubits l\'ogicos (codificados, con detecci\'on de errores) frente a qubits f\'isicos, sobre el
hardware Quantinuum H2. La idea cierra un hueco concreto en la literatura: existe un art\'iculo de
Nature (arXiv:2503.20870)~\cite{haghshenas2026} que calcula esos observables con qubits f\'isicos y publica sus datos
crudos, y existen trabajos que simulan con qubits l\'ogicos pero no calculan esos observables~\cite{dasu2026,froland2026}. La
comparaci\'on se hace contra los datos f\'isicos publicados de ese mismo experimento, con el encodeo
como \'unico eje que var\'ia. En orden de prioridad: (1)~l\'ogico vs f\'isico sobre observables f\'isicos,
(2)~frustraci\'on en la red triangular, (3)~termalizaci\'on, junto a hilos menores como la ley de
Gauss emergente para post-selecci\'on por simetr\'ia. Los observables locales de este tipo pueden simularse cl\'asicamente hasta cierto tama\~no~\cite{begusic2024}, por eso el proyecto se plantea como f\'isica m\'as metodolog\'ia de correcci\'on de errores y no como ventaja cu\'antica.

\subsection{Resumen en una p\'agina}
\begin{itemize}[leftmargin=1.4em]
  \item \textbf{La ruta a H2 funciona en las capas gratuitas.} El verificador de sintaxis de Quantinuum
    (H2-1SC) acept\'o y cotiz\'o \textbf{121 circuitos} sin errores. En el emulador sin ruido alojado en
    Nexus (H2-1LE) corrieron circuitos reales de extremo a extremo (subir, ejecutar, bajar, decodificar)
    con \textbf{costo reportado 0}, y los resultados coinciden con el valor exacto.
  \item \textbf{El precio es predecible.} La f\'ormula de costos que escribimos reproduce las cotizaciones
    del servidor \emph{al centavo}. Las compuertas $R_z$ (incluso condicionales) no se cobran.
  \item \textbf{Presupuesto.} El plan completo cuesta \textbf{39\,853 HQC de los 50\,000}
    (secci\'on~\ref{sec:presupuesto}); sobran unos 10\,100 para repeticiones.
  \item \textbf{Cambio de plan por el tama\~no del emulador.} El emulador H2 llega a 32 qubits, as\'i que
    el reticulado de Q1 pasa de $6{\times}8$ a $4{\times}5$ y el brazo $6{\times}6$ de escalado de ruido se
    reduce a $3{\times}4$ y $3{\times}6$ (decisi\'on m\'ia, ver secci\'on~\ref{sec:decisiones}).
  \item \textbf{Para encender el emulador} hay que editar \texttt{config.py} y pasar una aprobaci\'on
    escrita en la l\'inea de comandos. Una edici\'on de la configuraci\'on sola nunca gasta cr\'editos.
  \item \textbf{Revisi\'on independiente.} Dos revisores a ciegas encontraron 1 error mayor (ya
    corregido, con prueba de regresi\'on) y varios menores; ver la secci\'on~\ref{sec:revision}.
\end{itemize}

\subsection{C\'omo leer este reporte}
Las secciones 2 a 6 son cortas: qu\'e se construy\'o, qu\'e se logr\'o, el presupuesto, c\'omo encender el
emulador y qu\'e queda. \textbf{Los resultados f\'isicos con sus tablas y figuras est\'an en el
ap\'endice~\ref{ap:resultados}}; despu\'es vienen los problemas encontrados en el software de Quantinuum,
las fuentes y un registro de los identificadores internos (Q, K, FLAG, T, HW). Entre corchetes se
indica cu\'anto confiar en cada afirmaci\'on: \emph{establecido} (fuente verificada), \emph{derivado}
(demostrado y testeado aqu\'i), \emph{medido} (salida de nuestro c\'odigo o del servidor),
\emph{inferido} (razonado, no testeado). \textbf{No se corri\'o ning\'un experimento en el emulador con
ruido ni en hardware del H2}; los \'unicos n\'umeros de H2 son los datos \emph{publicados} de 2503.20870
y las \emph{cotizaciones} (precios) del servidor.

% ============================================================================
\section{Qu\'e se construy\'o}

Todo el c\'odigo est\'a detr\'as de un \'unico \texttt{config.py}. La base del proyecto es el n\'ucleo del TFIM,
los cuatro c\'odigos de correcci\'on, el TFIM codificado, la frustraci\'on, la termalizaci\'on, la ley de Gauss y la
comparaci\'on con los datos de Zenodo. Componentes para H2:

\begin{itemize}[leftmargin=1.4em]
  \item \textbf{Circuitos para H2} (\code{src/h2_circuits.py}, \code{src/h2_q3.py}). Cada programa
    (Iceberg~\cite{self2024}, el circuito sin codificar, Steane, 5 qubits, Iceberg concatenado~\cite{knill2005} y ahora la red triangular
    y la escalera de ruido) se escribe como \emph{un solo circuito} de pytket~\cite{sivarajah2020}, con medici\'on y reinicio
    a mitad de circuito. H2 y sus emuladores aceptan pytket, no Guppy; el descarte por s\'indrome se hace
    despu\'es, en el procesamiento cl\'asico.
  \item \textbf{Ruta de env\'io} (\code{src/h2_backend.py}, \code{src/h2_hosted.py}). Cuatro capas, de la m\'as
    barata a la m\'as cara: emulador local sin conexi\'on (gratis), verificador de sintaxis H2-1SC (gratis, da
    el precio), emulador sin ruido alojado en Nexus H2-1LE (gratis, $\le$20 qubits), y el emulador H2-1E
    (gasta HQC). Tres de las cuatro est\'an probadas en vivo.
  \item \textbf{Un ``driver'' de corridas} (\code{src/h2_driver.py}, \code{scripts/run_h2_nexus.py}).
    Convierte un grupo de corridas (G1, Q1, Q4, Q3\ldots) en circuitos compilados con su costo estimado y los
    ejecuta en la capa que diga \texttt{config.py}. Por defecto es un \emph{ensayo en seco}: no env\'ia nada.
  \item \textbf{Circuitos de Q3 para H2}: red triangular sin codificar y con Iceberg, y la escalera de
    escalado de ruido ($\lambda=1,2,3$ por inserci\'on de Paulis~\cite{haghshenas2026}, m\'as un punto de plegado de compuertas~\cite{giurgica2020}),
    todos compilados sin optimizaci\'on del servidor (si no, el plegado se cancelar\'ia en silencio).
  \item \textbf{Endurecimiento de la ruta de gasto}: la estimaci\'on de HQC ahora no cobra $R_z$
    (coincide con el servidor); las compuertas condicionales se cuentan por la compuerta que envuelven;
    el emulador alojado se manda con \texttt{no\_opt}; y l\'imites duros (32 qubits, 10\,000 shots por
    trabajo, 50 HQC por shot~\cite{quantinuumemul,quantinuumh2}).
  \item \textbf{Compilaci\'on aleatoria (RC)}: los circuitos codificados pueden ``torcerse'' con Paulis aleatorios (\emph{twirling})~\cite{wallman2016}, como hace el art\'iculo de referencia~\cite{haghshenas2026}. En la ruta a H2 se usa con Iceberg, salvo en Q3 (ver Decisiones).
  \item \textbf{Documentos}: \code{docs/nexus_toggle.md} (la ruta completa y c\'omo encender el emulador),
    \code{docs/h2_budget.md} (presupuesto) y \code{docs/quantinuum_issues.md} (el registro de problemas de Guppy/Selene, qnexus,
    pytket-quantinuum, PECOS y TKET~\cite{qilog}).
\end{itemize}

Las pruebas: \textbf{662 pruebas (m\'as 23 omitidas a prop\'osito)} pasan en el entorno principal y \textbf{288} en el entorno del emulador local
(dos entornos porque el paquete del emulador local no convive con el de Selene). Los archivos de pruebas
de la ruta a H2 se escribieron \emph{antes} que el c\'odigo y por un agente que no ve la implementaci\'on; sus huellas
(SHA-256) est\'an en \code{handoff/14-h2-finish.md}.

% ============================================================================
\section{Qu\'e se logr\'o}

\subsection{Capas gratuitas probadas en vivo en Nexus}
\begin{table}[!ht]\centering\small
\caption{Qu\'e se prob\'o en las capas gratuitas (17 de octubre). Todo es \emph{plomer\'ia y precio}, no
f\'isica con ruido. Fuentes: \code{data/h2\_sc\_quotes*.txt}, \code{data/h2\_hosted\_check.txt},
\code{data/h2\_hosted\_q3.txt}.}
\begin{tabular}{p{4.3cm}p{3.1cm}p{6.6cm}}
\toprule
prueba & resultado & qu\'e demuestra\\
\midrule
verificador de sintaxis H2-1SC & 121 circuitos aceptados & el servidor acepta nuestros circuitos, incluida la medici\'on a mitad de circuito y las compuertas condicionales\\
precio ``al centavo'' & estimaci\'on = cotizaci\'on & nuestra f\'ormula de HQC es correcta; $R_z$ no se cobra\\
emulador alojado H2-1LE (primera tanda) & 8 circuitos, costo 0.0 & subir, ejecutar, bajar y decodificar funciona; $Z_2$ dentro de 1.5$\sigma$ de lo exacto\\
emulador alojado H2-1LE (Q3) & 79 circuitos, costo 0.0 & los circuitos de Q3 (red triangular, Iceberg, escalera) dan lo mismo que una evoluci\'on independiente en NumPy: 19 puntos $\times$ 3 observables, peor $|z|=2.52$; Iceberg con aceptaci\'on 1.000\\
\bottomrule
\end{tabular}
\end{table}

\subsection{Una correcci\'on de precio}
Mi primera estimaci\'on cobraba las compuertas $R_z$ y daba un 15--20\,\% de m\'as. El servidor demostr\'o que no
se cobran. Tambi\'en descubr\'i un error m\'io: las compuertas \emph{condicionales} se contaban mal, lo que
dejaba una diferencia de hasta 0.5\,\% en Steane y Iceberg concatenado. Corregido, la estimaci\'on coincide
con el servidor al centavo en los 9 circuitos de Q4 [medido].

% ============================================================================
\section{Presupuesto HQC}\label{sec:presupuesto}

Los 50\,000 HQC de la p\'agina de uso son cr\'editos del \emph{emulador} H2-1E/2E~\cite{quantinuumemul}; no hay cr\'editos de hardware
real. Todos los precios de la tabla son \textbf{cotizaciones del servidor} (verificador gratuito), con 1000
shots por punto en G1, Q1 y Q4 y 500 en Q3 [medido]. Aparte, la estimaci\'on sin conexi\'on coincide con cada
cotizaci\'on al centavo.

\begin{table}[h]\centering\small
\caption{Costo del plan completo en HQC (cotizado). \emph{Nada de esto se ha corrido en el emulador.}}
\begin{tabular}{llrrr}
\toprule
grupo & qu\'e es & qubits & circuitos & HQC\\
\midrule
G1 & Iceberg $2{\times}2$ (simple y RC) + sin codificar & 8 / 4 & 66 & 1\,901\\
Q1 & Iceberg $4{\times}5$ (RC) + sin codificar $4{\times}5$, pasos 1--6 & 24 / 20 & 126 & 11\,636\\
Q4 & Steane, 5 qubits, concatenado, pasos 1--3 & 16 / 12 / 18 & 9 & 9\,817\\
Q3-bare & triangular $3{\times}4$ sin codificar, $t=1$--8 & 12 & 8 & 5\,670\\
Q3-iceberg & triangular $3{\times}4$ Iceberg, $t=1$--4 & 16 & 4 & 5\,283\\
Q3-ladder & escalera de ruido $3{\times}4$ ($\lambda=1,2,3$ + plegado) & 12 & 64 & 4\,093\\
Q3-bare-3x6 & triangular $3{\times}6$ sin codificar, $t=1$--3 & 18 & 3 & 1\,453\\
\midrule
\textbf{Total} & & & & \textbf{39\,853}\\
\textbf{Reserva} & de los 50\,000 & & & \textbf{10\,147}\\
\bottomrule
\end{tabular}
\end{table}

La estimaci\'on vieja de Q3 (unos 10\,000) era una conjetura por f\'ormula y qued\'o 6\,500 corta: la curva sin
codificar llega a $t=8$ (32 pasos de Trotter) e Iceberg a $t=4$. Si hace falta recortar: sin codificar
hasta $t\le6$ ahorra 2\,339, Iceberg hasta $t\le3$ ahorra 2\,012, o la mitad de los shots en Q3 ahorra casi la
mitad. Es decisi\'on de ustedes (secci\'on~\ref{sec:decisiones}).

La corrida m\'as barata de punta a punta es el circuito sin codificar $2{\times}2$ a un paso: \textbf{23.40 HQC}
con 1000 shots. Ah\'i empezar\'ia. También dejé estas dimensiones de grillas porque muestran comportamiento interesante y están dentro de los HQCs que tenemos cada uno, pero perfectamente podemos simular Q3-bare con 6x6 en lugar de 3x6, solo que eso tomaría mucho más HQCs.

% ============================================================================
\section{C\'omo encender el emulador}\label{sec:encender}

Es lo \'unico que falta, y lo dej\'e a prop\'osito como un cambio de configuraci\'on (de \texttt{docs/nexus\_toggle.md}):
\begin{enumerate}[leftmargin=1.6em]
  \item En \texttt{config.py}, \texttt{HardwareConfig}: \texttt{nexus\_tier = "emulator"},
    \texttt{allow\_hardware = True}, \texttt{max\_cost\_hqc = <tope en HQC>} y el nombre del proyecto de Nexus.
  \item Iniciar sesi\'on en Nexus (c\'odigo de dispositivo que se aprueba en el navegador); es por
    contenedor.
  \item Correr el driver con la \textbf{aprobaci\'on por corrida}, que ninguna edici\'on de la configuraci\'on
    reemplaza: \texttt{python -I scripts/run\_h2\_nexus.py --group bare-2x2-step1 --approval "<frase>"}.
\end{enumerate}
Antes de gastar, el driver construye los circuitos, \textbf{vuelve a cotizarlos en el verificador gratuito} y se
niega a enviar si el precio supera el tope. Tambi\'en se niega con m\'as de 32 qubits, m\'as de 10\,000 shots,
m\'as de 50 HQC por shot~\cite{quantinuumh2}, o con cualquier dispositivo que no sea un emulador (el hardware real queda
excluido). La l\'inea de comandos no puede elegir el emulador por s\'i sola. Orden sugerido, cada paso con mi
visto bueno: sin codificar $2{\times}2$ a un paso $\to$ G1 $\to$ comparar con el emulador local, Selene y Stim
$\to$ Q1 $\to$ Q4 $\to$ Q3.

% ============================================================================
\section{Qu\'e queda por hacer}\label{sec:queda}

\paragraph{Depende de gastar cr\'editos (cada paso con aprobaci\'on expl\'icita).}
\begin{enumerate}[leftmargin=1.6em]
  \item \textbf{Calibraci\'on (HW-5).} Primero el circuito m\'as chico en el emulador H2-1E, comparado con las
    predicciones locales antes de gastar m\'as.
  \item \textbf{Medir el ruido} (fracci\'on ZZ y sobre-rotaci\'on del gate de dos qubits) y volver a correr la
    tabla de sensibilidad con esos valores.
  \item \textbf{La prueba decisiva (Q1).} Codificado vs desnudo con los mismos shots: ?`el encodeo ayuda?
  \item \textbf{Escalado de ruido en Q3} ($3{\times}4$) y los dem\'as grupos del plan.
\end{enumerate}


% ============================================================================
\section{Decisiones}\label{sec:decisiones}

\paragraph{Ya decididas.}
\begin{itemize}[leftmargin=1.4em]
  \item \textbf{Reticulado de Q1: $4{\times}5$} (24 qubits codificado, 20 sin codificar), en lugar de $6{\times}8$
    (52 qubits, no cabe en el emulador de 32~\cite{quantinuumemul}). Tambi\'en se puede cruzar con el simulador local.
  \item \textbf{El brazo $6{\times}6$ se reduce a $3{\times}4$ y $3{\times}6$}: la escalera de ruido en $3{\times}4$ y
    el tama\~no en $3{\times}6$. El $6{\times}6$ queda como dise\~no solo en Stim; no se corre en ning\'un lado.
  \item \textbf{La estimaci\'on de HQC no cobra $R_z$} por defecto (D-HQC); se cambi\'o una aserci\'on de las pruebas
    congeladas y se registr\'o la nueva huella.
\end{itemize}

\paragraph{Abiertas (necesito su opini\'on).}
\begin{enumerate}[leftmargin=1.6em]
  \item \textbf{Asignaci\'on de los 50\,000 HQC y cu\'anto recortar Q3.} El plan completo usa 39\,853.
  \item \textbf{Tope de costo por trabajo (D-MAXCOST).} Cuando un grupo se parte en varios trabajos (por
    ejemplo G1 y la escalera), cada trabajo recibe como tope el tope \emph{total}, as\'i que el seguro del
    servidor vale hasta $n\times$ el tope. El driver ya revisa antes que la suma de cotizaciones no pase el
    tope, y el precio real coincide con la cotizaci\'on; cambiar a un tope por trabajo exige modificar una prueba
    congelada, por eso lo dejo a su criterio.
  \item \textbf{?`Iceberg en Q3 con compilaci\'on aleatoria (RC)?} La decisi\'on D-RC dice que las corridas
    codificadas en H2 usan RC~\cite{wallman2016}, pero el circuito Iceberg de Q3 se construy\'o \emph{sin} RC para coincidir con
    la predicci\'on cl\'asica sin twirling. Recomiendo dejarlo as\'i para la primera pasada y agregar RC si la
    calibraci\'on lo pide.
  \item \textbf{La frase de aprobaci\'on es una constante del repositorio}, as\'i que evita gastos accidentales o
    en bucle pero no a alguien que la copie a prop\'osito. Me parece aceptable; d\'iganme si prefieren algo m\'as
    estricto.
\end{enumerate}

\paragraph{Pendientes estacionadas.} Protecci\'on de gauge activa (K1), orden de 3 subredes (K3), difusi\'on (K5) y
Z2 vs U(1) (K6): ninguna construida, todas opcionales.

% ============================================================================
\section{Revisi\'on independiente}\label{sec:revision}
Al final de la ronda hice una sola pasada con dos revisores a ciegas (modelo Opus, sin ver mi razonamiento):
uno sobre la seguridad del gasto y otro sobre los circuitos, los n\'umeros y el precio. Resultado en
\code{review/17-h2-finish.md}:
\begin{itemize}[leftmargin=1.4em]
  \item \textbf{1 error mayor, corregido.} Un argumento expl\'icito \texttt{tier="emulator"} pod\'ia saltarse la
    configuraci\'on (la l\'inea de comandos no, pero la funci\'on s\'i). Ahora la configuraci\'on tiene que decir
    ``emulator''. Prueba de regresi\'on a\~nadida; fall\'o antes y pasa ahora.
  \item \textbf{Menores corregidos}: cotizaciones incompletas, l\'imites NaN o infinitos que apagaban una
    compuerta, tres huecos de pruebas (m\'ascara de bobinado, pasos de la escalera, semillas de inserci\'on) y un
    mensaje de error confuso.
  \item \textbf{Confirmado por los revisores}: la f\'isica de los circuitos (reconstruida por otro m\'etodo, con
    solapamiento 1.0), los n\'umeros de los archivos de datos recalculados a mano, y el precio al centavo.
\end{itemize}
Ning\'un camino gasta cr\'editos sin la configuraci\'on correcta, la aprobaci\'on por corrida y un tope finito.

% ============================================================================
\appendix

\section{Resultados f\'isicos}\label{ap:resultados}

Son resultados cl\'asicos, sin ruido de hardware. El brazo de escalado de ruido ya no es $6{\times}6$ (secci\'on~\ref{sec:decisiones}).

\subsection{Q1 --- L\'ogico vs f\'isico}

\paragraph{Chequeo contra el art\'iculo~\cite{haghshenas2026}.} Los 112 pares de dos qubits por paso coinciden exactamente con un toro
cuadrado peri\'odico $7\times8$ con numeraci\'on por filas [medido], y reproducir los circuitos compuerta por
compuerta da lo mismo que nuestra evoluci\'on con $(J,h)=(+1,+2)$, $dt=0.25$ [derivado, testeado].

\begin{table}[h]\centering\small
\caption{Quench 1 ($\theta=30^\circ$): correlaci\'on media entre pares. ``H2 publicado'' son los datos
de 56 qubits \emph{del art\'iculo 2503.20870}~\cite{haghshenas2026,zenodo}, no nuestros. ``Nuestra sim'' es una simulaci\'on sin ruido
en $4\times5$. Fuente: \code{data/q1\_compare.txt}.}
\begin{tabular}{rccc}
\toprule
paso & nuestra sim ($4\times5$) & H2 publicado, mitigado & H2 publicado, crudo\\
\midrule
1 & 0.761 & $0.732\pm0.009$ & $0.726\pm0.007$\\
4 & 0.695 & $0.623\pm0.017$ & $0.582\pm0.013$\\
20 & 0.690 & $0.554\pm0.040$ & $0.379\pm0.016$\\
\bottomrule
\end{tabular}
\end{table}
\noindent Muestra que nuestro montaje reproduce los circuitos publicados. No dice nada de qubits
l\'ogicos: nuestras redes chicas difieren de 56 qubits por efectos de tama\~no de hasta 0.06--0.33.

\paragraph{?`El encodeo ayuda?} Un error de dos qubits de tipo $XX$, $YY$ o $ZZ$ sobre el par en
que act\'ua una compuerta l\'ogica es invisible para Iceberg: 3 de los 15 errores posibles [derivado].
Si una fracci\'on $\eta$ del error es $ZZ$ extra, la parte invisible crece.

\begin{table}[h]\centering\small
\caption{Error no detectado por paso relativo al circuito desnudo, en la red de prueba $2\times2$,
con el prior del H2 ($p_2=6\times10^{-4}$, \cite{moses2023,haghshenas2026}) \pred. Menor que 1 significa que el encodeo ayuda. En
redes m\'as grandes los cocientes son menores (Iceberg 0.44 en cuadradas grandes, 0.375 en la
triangular, con $\eta=0.05$). Fuente: \code{data/noise\_sensitivity.txt}.}
\begin{tabular}{lcccc}
\toprule
escenario & Iceberg & Steane & 5 qubits & concat.\ $d=4$\\
\midrule
prior H2, $\eta=0.05$ & $\mathbf{0.60}$ & 0.53 & 0.52 & 0.53\\
$\eta$ entre 0.03 y 0.08 & 0.56--0.66 & 0.48--0.59 & 0.48--0.59 & 0.48--0.59\\
m\'as sobre-rotaci\'on $\sin^2\epsilon=10^{-3}$ & 0.92 & 0.91 & 0.91 & 0.91\\
prueba de estr\'es $\eta=0.5$ ($10\times$ el prior) & 1.50 & 1.50 & 1.50 & 1.50\\
\midrule
shots que sobreviven tras 3 pasos & 0.92 & 0.57 & 0.64 & 0.64\\
\bottomrule
\end{tabular}
\end{table}
\noindent El cruce en $\eta\approx0.25$ vale para la red de prueba $2\times2$ (peor caso); en redes m\'as grandes
sube a $\eta\approx0.36$--$0.55$, y el prior de la literatura ($\eta\approx0.04$--$0.05$) queda a\'un m\'as lejos.
El $\eta$ del prior sale de leer un gr\'afico publicado~\cite{haghshenas2026} [inferido] y la sobre-rotaci\'on no est\'a publicada,
as\'i que el \emph{go} es una predicci\'on; solo el emulador o la m\'aquina lo confirman.

\subsection{Q4 --- Qu\'e c\'odigo y por qu\'e}
\begin{itemize}[leftmargin=1.4em]
  \item Los circuitos de medici\'on de s\'indrome de los cuatro c\'odigos son tolerantes a fallos: ning\'un
    fallo \'unico ah\'i produce un error l\'ogico sin detectar (verificado con Stim~\cite{gidney2021}; s\'indrome extra\'ido con una sola qubit auxiliar y una bandera~\cite{chao2018}) [medido]. Todo el error indetectable viene
    de las compuertas de rotaci\'on l\'ogicas, sin importar la distancia [derivado y medido].
  \item Por eso los c\'odigos de distancia 3 y 4 solo ganan alrededor de $1.2\times$ en error no
    detectado, a cambio de conservar 3--4 veces menos shots con $p_2=10^{-3}$ y de usar 3--7 veces m\'as
    compuertas de dos qubits \pred.
  \item \textbf{Iceberg~\cite{self2024} es el c\'odigo recomendado.} Los otros est\'an construidos y funcionan, solo que
    quedan dominados. Selene conserva menos shots que Stim porque cobra ruido en cada operaci\'on nativa de
    un qubit; es una diferencia de definici\'on, no un bug.
\end{itemize}

\subsection{Q2 --- Frustraci\'on: triangular vs cuadrada}
El mapeo de la red triangular antiferromagn\'etica a d\'imeros~\cite{moessner2001} y su diagrama de fases~\cite{isakov2003} dan el marco. Estados base exactos, sin ruido ni error de muestreo [medido].

\begin{table}[h]\centering\small
\caption{Correlaciones en el par $4\times4$, $J=-1$. $C(1)$ = vecinos m\'as cercanos; ``lejos'' es la
mayor distancia de cada toro. Fuente: \code{data/frustration\_correlations.txt}.}
\begin{tabular}{ccc|cc}
\toprule
campo $h$ & cuadrada $C(1)$ & cuadrada, lejos & triangular $C(1)$ & triangular, lejos\\
\midrule
0.1 & $-0.999$ & $+0.999$ & $-0.332$ & $+0.442$\\
0.3 & $-0.994$ & $+0.994$ & $-0.325$ & $+0.413$\\
0.6 & $-0.977$ & $+0.977$ & $-0.268$ & $+0.196$\\
1.0 & $-0.936$ & $+0.936$ & $-0.204$ & $+0.068$\\
\bottomrule
\end{tabular}
\end{table}
\begin{itemize}[leftmargin=1.4em]
  \item \textbf{La firma de frustraci\'on:} la correlaci\'on entre vecinos en la triangular es 3 a 5 veces
    m\'as d\'ebil, cercana a $-1/3$. Un toro cuadrado de largo impar est\'a frustrado ($3\times4$ no era
    una referencia limpia); $4\times4$ es el \'unico toro cuadrado limpio al alcance de la diagonalizaci\'on exacta.
  \item El patr\'on triangular a larga distancia cambia de signo seg\'un la forma del toro, as\'i que no
    contamos una historia sobre eso, y no afirmamos orden de reloj (K3 no est\'a construido).
\end{itemize}

\subsection{Q3 --- Termalizaci\'on}
\paragraph{La referencia.} La termalizaci\'on ocurre solo dentro de cada sector de bobinado~\cite{lan2017}, as\'i que la
referencia es por sector (en $3\times4$: densidad de flipabilidad 0.267 en el sector vs 0.235 global) [medido].
En $3\times6$ la densidad baja de 0.667 a unos 0.383, el valor de tiempo largo del c\'umulo finito (0.388), no el
de temperatura infinita del sector (0.331); la diferencia se achica con el tama\~no [inferido].

\begin{figure}[h]\centering
\figura{q3_relaxation.png}{0.9\linewidth}
\caption{Relajaci\'on ideal en $3\times6$ a dos campos. Simulaci\'on cl\'asica, sin ruido.}
\end{figure}

\paragraph{?`El paso de tiempo es suficientemente chico?} Existe un umbral n\'itido del paso de Trotter por debajo del cual los observables locales se mantienen fieles~\cite{heyl2019}, e incluso un Trotter grueso pre-termaliza hacia el valor t\'ermico del hamiltoniano objetivo~\cite{yang2023}. En nuestro caso el paso deja de ser fiel cerca de
$dt\approx\pi/6=0.52$. Nuestro $dt=0.25$ queda 1.3 a 2 veces por debajo con campos $h\le0.3$ [medido].

\begin{figure}[h]\centering
\figura{trotter_threshold.png}{0.8\linewidth}
\caption{Error de Trotter vs paso de tiempo. Cl\'asico, sin ruido.}
\end{figure}

\paragraph{Con ruido \pred.}
\begin{itemize}[leftmargin=1.4em]
  \item Sin mitigar, la curva desnuda se calienta. \textbf{Post-seleccionar con la ley de Gauss reduce el error
    unas 30 veces temprano y 21 tarde, con 1.4 a 4.8 veces m\'as shots.} Iceberg solo da alrededor de
    $2.7\times$ con 3.1 veces m\'as shots, y sobrevive hasta $t\approx12$.
  \item El ruido empuja las curvas hacia el valor de temperatura infinita: en hardware \textbf{el ruido puede
    imitar la termalizaci\'on}; por eso la escalera de ruido ($\lambda=1,2,3$), por inserci\'on de Paulis~\cite{haghshenas2026} y por plegado de compuertas~\cite{giurgica2020}.
\end{itemize}
\begin{figure}[h]\centering
\figura{q3_noisy.png}{0.85\linewidth}
\caption{$3\times4$: ideal vs desnudo, con Gauss, Iceberg e Iceberg con Gauss. Predicci\'on cl\'asica con ruido asumido; no es H2.}
\end{figure}

\subsection{Ley de Gauss emergente como mitigaci\'on}
La post-selecci\'on por ley de Gauss ya se ha usado en hardware con generadores locales~\cite{mildenberger2025}. Veredicto con ruido de Pauli por qubit asumido, en $3\times3$ y $3\times4$: \textbf{no sirve para correcci\'on}
(distancia 1); \textbf{sirve modestamente para detecci\'on solo contra ruido $X/Y$} (2 a 8 veces menos sesgo con 1.1 a
1.8 veces m\'as shots, con un piso de orden $(h/J)^2$); no sirve contra defasaje; la automitigaci\'on no sirve como
m\'etodo general.

% ============================================================================
\section{Problemas encontrados en el software de Quantinuum}
Como tengo acceso a Nexus tambi\'en para ayudar a depurar su software, cada problema queda en
\code{docs/quantinuum\_issues.md} con su reproducci\'on, listo para enviar. Ninguno se ha enviado todav\'ia. Cada entrada del registro~\cite{qilog} trae las versiones, la reproducci\'on y qu\'e pedirle a Quantinuum.
\begin{table}[h]\centering\small
\begin{tabular}{lp{11cm}}
\toprule
id & resumen\\
\midrule
QI-1 & el emulador local (\texttt{pytket-quantinuum[pecos]}) no convive con la versi\'on de Selene que usamos: dos entornos\\
QI-2 & el emulador local ignora \texttt{noisy\_simulation} y es siempre sin ruido\\
QI-3 & el cliente HUGR (Guppy) de Nexus est\'a tipado para Helios; el soporte para H2 es ``limitado''. Por eso la ruta a H2 usa un circuito pytket y no un programa Guppy\\
QI-4 & \texttt{qnexus.hugr.cost()} est\'a obsoleta (se usa \texttt{cost\_confidence()}); solo un aviso\\
QI-5 & \texttt{circuits.cost} arma el nombre ``H2-1ESC'' desde una configuraci\'on E (un \texttt{strip} descartado)\\
QI-6 & el compilador HUGR-QIS de Selene (Guppy) rechaza la compuerta \texttt{XXPhase} de pytket; se escribe como H$\cdot$ZZPhase$\cdot$H\\
QI-7 & el modelo despolarizante de Selene cobra el error de un qubit en cada operaci\'on nativa (rxy y rz), unas 3.85 veces m\'as exposici\'on que contando a nivel de pytket; baja la aceptaci\'on frente a Stim\\
QI-8 & la optimizaci\'on de TKET (niveles 1--2) cancela en silencio las compuertas plegadas\\
QI-9, 11 & PECOS reescala los errores y no cobra la inicializaci\'on inicial\\
QI-10 & pytket$\to$PHIR descarta bits cl\'asicos que ning\'un comando toca (ya conocido por Quantinuum, citado en su c\'odigo)\\
QI-12 & las p\'aginas de emuladores se contradicen (26 vs 20 qubits)\\
QI-13 & aviso de un endpoint ``que se borra el 30 nov 2025'', fecha ya pasada\\
QI-14 & \texttt{start\_execute\_job} se cae con circuitos crudos pese a un decorador con ese nombre\\
QI-15 & los l\'imites de HQC (50 por shot, 10\,000 shots) solo se documentan para ``hardware''; $R_z$ condicional no se cobra y no est\'a documentado\\
\bottomrule
\end{tabular}
\end{table}

% ============================================================================
\section{Andamios del setup}
Para que no cunda el p\'anico cuando vean el proyecto, especifico c\'omo us\'e la IA y qu\'e hay en cada carpeta:
\begin{enumerate}
  \item \texttt{SCOPE.md}: la ``constituci\'on'' del proyecto (las cuatro preguntas, los l\'imites de tama\~no, qu\'e se
    puede hacer sin preguntarme). Si hay duda, manda este documento; c\'ambienlo a consciencia.
  \item \texttt{CLAUDE.md}: las reglas de trabajo del agente (qu\'e est\'a prohibido sin mi aprobaci\'on, convenciones de
    f\'isica, trampas de las herramientas).
  \item \texttt{DECISIONS.md}: decisiones abiertas y tomadas (incluye la decisi\'on del reticulado de Q1 y del brazo de escalado de ruido, y D-HQC).
  \item \texttt{config.py}, \texttt{main.py}, \texttt{requirements*.txt}, \texttt{src/}: la configuraci\'on \'unica y el
    c\'odigo. Dos archivos de requisitos: el principal y el del emulador local (\texttt{requirements-le.txt}).
  \item \texttt{scripts/}: los scripts que generan cada resultado. Para H2: \texttt{run\_h2\_nexus.py} (el driver),
    \texttt{h2\_sc\_quotes*.py} (cotizaciones gratuitas), \texttt{h2\_hosted\_*.py} (emulador alojado).
  \item \texttt{data/}: resultados en texto plano; cada uno dice qu\'e script lo gener\'o y si es una predicci\'on cl\'asica con ruido asumido.
  \item \texttt{docs/}, \texttt{review/}, \texttt{research/}: la explicaci\'on de cada resultado, el registro de problemas del software de Quantinuum (\texttt{docs/quantinuum\_issues.md}), las revisiones
    independientes numeradas (la 17 cubre la ruta a H2) y las notas de investigaci\'on.
  \item \texttt{handoff/}: los encargos y reportes entre sesiones; \texttt{handoff/14-h2-finish.md} resume el trabajo de la ruta a H2.
  \item \texttt{tests/}: las pruebas. Las ``congeladas'' se escribieron antes del c\'odigo y est\'an fijadas por huella (SHA-256).
  \item \texttt{Guppy on Quantinuum Systems Application.pdf}: la aplicaci\'on original, para corroborar que lo hecho corresponde a lo que dijimos que har\'iamos.
  \item Todo esto est\'a en la rama \texttt{claude/awesome-lamport-wun6jh} del repositorio.
\end{enumerate}

% ============================================================================
\begin{thebibliography}{99}
\addcontentsline{toc}{section}{Referencias}

\bibitem{haghshenas2026} R.~Haghshenas, E.~Chertkov, M.~Mills, \emph{et al.} (59 autores), ``Digital quantum magnetism on a trapped-ion quantum computer'', Nature \textbf{653}, 56--62 (2026), doi:\nolinkurl{10.1038/s41586-026-10445-3}, arXiv:2503.20870.

\bibitem{zenodo} Datos publicados de~\cite{haghshenas2026}, Zenodo, doi:\nolinkurl{10.5281/zenodo.18487089}.

\bibitem{self2024} C.~N.~Self, M.~Benedetti y D.~Amaro, ``Protecting Expressive Circuits with a Quantum Error Detection Code'', Nat.\ Phys.\ \textbf{20}, 219--224 (2024), doi:\nolinkurl{10.1038/s41567-023-02282-2}, arXiv:2211.06703.

\bibitem{chao2018} R.~Chao y B.~W.~Reichardt, ``Quantum error correction with only two extra qubits'', Phys.\ Rev.\ Lett.\ \textbf{121}, 050502 (2018), doi:\nolinkurl{10.1103/PhysRevLett.121.050502}, arXiv:1705.02329.

\bibitem{knill2005} E.~Knill, ``Quantum computing with realistically noisy devices'', Nature \textbf{434}, 39--44 (2005), doi:\nolinkurl{10.1038/nature03350}, arXiv:quant-ph/0410199.

\bibitem{dasu2026} S.~Dasu, M.~DeCross, A.~Y.~Guo, \emph{et al.} (45 autores), ``Computing with many encoded logical qubits beyond break-even'', arXiv:2602.22211 (2026).

\bibitem{froland2026} H.~Froland, D.~M.~Grabowska, S.~Grieninger, \emph{et al.} (11 autores), ``Realizing Error Suppression in Partially Fault-Tolerant Quantum Simulations with IBM Quantum Computers'', arXiv:2607.24947 (2026).

\bibitem{moessner2001} R.~Moessner y S.~L.~Sondhi, ``Ising models of quantum frustration'', Phys.\ Rev.\ B \textbf{63}, 224401 (2001), doi:\nolinkurl{10.1103/PhysRevB.63.224401}, arXiv:cond-mat/0011250.

\bibitem{isakov2003} S.~V.~Isakov y R.~Moessner, ``Interplay of quantum and thermal fluctuations in a frustrated magnet'', Phys.\ Rev.\ B \textbf{68}, 104409 (2003), doi:\nolinkurl{10.1103/PhysRevB.68.104409}, arXiv:cond-mat/0302105.

\bibitem{lan2017} Z.~Lan y S.~Powell, ``Eigenstate thermalization hypothesis in quantum dimer models'', Phys.\ Rev.\ B \textbf{96}, 115140 (2017), doi:\nolinkurl{10.1103/PhysRevB.96.115140}, arXiv:1706.02601.

\bibitem{begusic2024} T.~Begu\v{s}i\'c, J.~Gray y G.~K.-L.~Chan, ``Fast and converged classical simulations of evidence for the utility of quantum computing before fault tolerance'', Sci.\ Adv.\ \textbf{10}, eadk4321 (2024), doi:\nolinkurl{10.1126/sciadv.adk4321}, arXiv:2308.05077.

\bibitem{heyl2019} M.~Heyl, P.~Hauke y P.~Zoller, ``Quantum localization bounds Trotter errors in digital quantum simulation'', Sci.\ Adv.\ \textbf{5}, eaau8342 (2019), doi:\nolinkurl{10.1126/sciadv.aau8342}, arXiv:1806.11123.

\bibitem{yang2023} Y.~Yang, A.~Christianen, S.~Coll-Vinent, \emph{et al.} (8 autores), ``Simulating prethermalization using near-term quantum computers'', PRX Quantum \textbf{4}, 030320 (2023), doi:\nolinkurl{10.1103/PRXQuantum.4.030320}, arXiv:2303.08461.

\bibitem{mildenberger2025} J.~Mildenberger, W.~Mruczkiewicz, J.~C.~Halimeh, Z.~Jiang y P.~Hauke, ``Probing confinement in a $Z_2$ lattice gauge theory on a quantum computer'', Nat.\ Phys.\ \textbf{21}, 312--317 (2025), doi:\nolinkurl{10.1038/s41567-024-02723-6}, arXiv:2203.08905.

\bibitem{moses2023} S.~A.~Moses, C.~H.~Baldwin, M.~S.~Allman, \emph{et al.} (96 autores), ``A Race Track Trapped-Ion Quantum Processor'', Phys.\ Rev.\ X \textbf{13}, 041052 (2023), doi:\nolinkurl{10.1103/PhysRevX.13.041052}, arXiv:2305.03828.

\bibitem{giurgica2020} T.~Giurgica-Tiron, Y.~Hindy, R.~LaRose, A.~Mari y W.~J.~Zeng, ``Digital zero noise extrapolation for quantum error mitigation'', 2020 IEEE International Conference on Quantum Computing and Engineering (QCE), doi:\nolinkurl{10.1109/QCE49297.2020.00045}, arXiv:2005.10921.

\bibitem{wallman2016} J.~J.~Wallman y J.~Emerson, ``Noise tailoring for scalable quantum computation via randomized compiling'', Phys.\ Rev.\ A \textbf{94}, 052325 (2016), doi:\nolinkurl{10.1103/PhysRevA.94.052325}, arXiv:1512.01098.

\bibitem{gidney2021} C.~Gidney, ``Stim: a fast stabilizer circuit simulator'', Quantum \textbf{5}, 497 (2021), doi:\nolinkurl{10.22331/q-2021-07-06-497}, arXiv:2103.02202.

\bibitem{sivarajah2020} S.~Sivarajah, S.~Dilkes, A.~Cowtan, W.~Simmons, A.~Edgington y R.~Duncan, ``t$|$ket$\rangle$: A Retargetable Compiler for NISQ Devices'', doi:\nolinkurl{10.1088/2058-9565/ab8e92}, arXiv:2003.10611.

\bibitem{quantinuumemul} Quantinuum, ``Quantinuum Emulators'', documentaci\'on de Quantinuum Systems, \url{https://docs.quantinuum.com/systems/trainings/h2/getting_started/emulators.html} (le\'ida el 17 de octubre de 2026).

\bibitem{quantinuumh2} Quantinuum, ``H2 Operation'', documentaci\'on de Quantinuum Systems, \url{https://docs.quantinuum.com/systems/user_guide/hardware_user_guide/h2.html} (le\'ida el 17 de octubre de 2026).

\bibitem{qilog} Registro de problemas del software de Quantinuum (QI-1 a QI-15), \texttt{docs/quantinuum\_issues.md}, repositorio KmarH/TFIM\_QEC\_Frustration, rama \texttt{claude/awesome-lamport-wun6jh}, 17 de octubre de 2026. Versiones: pytket 2.18.4, pytket-quantinuum 0.59.3, qnexus 0.51.1, guppylang 1.1.2, selene-sim 0.3.4 (0.2.18 en el entorno del emulador local).

\end{thebibliography}
\end{document}
